\section{AHSME 1988}
        \begin{enumerate}
        	\item (\textit{AHSME 1988}) $\sqrt{8}+\sqrt{18}=$
$\text{(A)}\ \sqrt{20} \qquad  \text{(B)}\ 2(\sqrt{2}+\sqrt{3}) \qquad  \text{(C)}\ 7 \qquad  \text{(D)}\ 5\sqrt{2}\qquad \text{(E)}\ 2\sqrt{13}$


	\item (\textit{AHSME 1988}) Triangles $ABC$ and $XYZ$ are similar, with $A$ corresponding to $X$ and $B$ to $Y$. If $AB=3, BC=4$, and $XY=5$, then $YZ$ is: 


 $\text{(A)}\ 3\frac{3}{4} \qquad  \text{(B)}\ 6 \qquad  \text{(C)}\ 6\frac{1}{4} \qquad  \text{(D)}\ 6\frac{2}{3} \qquad  \text{(E)}\ 8$


	\item (\textit{AHSME 1988}) 
                    \begin{figure}[H]
                        \centering
                        \includegraphics[width=0.5\textwidth]{Problem/AHSME_I_1988_3_Problem_Figure_1.png}
                    \end{figure}
                    


 Four rectangular paper strips of length $10$ and width $1$ are put flat on a table and overlap perpendicularly as shown. How much area of the table is covered?


 $\text{(A)}\ 36 \qquad  \text{(B)}\ 40 \qquad  \text{(C)}\ 44 \qquad  \text{(D)}\ 98 \qquad  \text{(E)}\ 100$


	\item (\textit{AHSME 1988}) The slope of the line $\frac{x}{3} + \frac{y}{2} = 1$ is


 $\textbf{(A)}\ -\frac{3}{2}\qquad \textbf{(B)}\ -\frac{2}{3}\qquad \textbf{(C)}\ \frac{1}{3}\qquad \textbf{(D)}\ \frac{2}{3}\qquad \textbf{(E)}\ \frac{3}{2}$


	\item (\textit{AHSME 1988}) If $b$ and $c$ are constants and $(x + 2)(x + b) = x^2 + cx + 6$, then $c$ is


 $\textbf{(A)}\ -5\qquad \textbf{(B)}\ -3\qquad \textbf{(C)}\ -1\qquad \textbf{(D)}\ 3\qquad \textbf{(E)}\ 5$


	\item (\textit{AHSME 1988}) A figure is an equiangular parallelogram if and only if it is a


 $\textbf{(A)}\ \text{rectangle}\qquad \textbf{(B)}\ \text{regular polygon}\qquad \textbf{(C)}\ \text{rhombus}\qquad \textbf{(D)}\ \text{square}\qquad \textbf{(E)}\ \text{trapezoid}$


	\item (\textit{AHSME 1988}) Estimate the time it takes to send $60$ blocks of data over a communications channel if each block consists of $512$
"chunks" and the channel can transmit $120$ chunks per second.


 $\textbf{(A)}\ 0.04 \text{ seconds}\qquad \textbf{(B)}\ 0.4 \text{ seconds}\qquad \textbf{(C)}\ 4 \text{ seconds}\qquad \textbf{(D)}\ 4\text{ minutes}\qquad \textbf{(E)}\ 4\text{ hours}$


	\item (\textit{AHSME 1988}) If $\frac{b}{a} = 2$ and $\frac{c}{b} = 3$, what is the ratio of $a + b$ to $b + c$?


 $\textbf{(A)}\ \frac{1}{3}\qquad \textbf{(B)}\ \frac{3}{8}\qquad \textbf{(C)}\ \frac{3}{5}\qquad \textbf{(D)}\ \frac{2}{3}\qquad \textbf{(E)}\ \frac{3}{4}$


	\item (\textit{AHSME 1988}) 
                    \begin{figure}[H]
                        \centering
                        \includegraphics[width=0.5\textwidth]{Problem/AHSME_I_1988_9_Problem_Figure_1.png}
                    \end{figure}
                    


 An $8' \times 10'$ table sits in the corner of a square room, as in Figure $1$ below. 
The owners desire to move the table to the position shown in Figure $2$. 
The side of the room is $S$ feet. What is the smallest integer value of $S$ for which the table can be moved as desired without tilting it or taking it apart?


 $\textbf{(A)}\ 11\qquad \textbf{(B)}\ 12\qquad \textbf{(C)}\ 13\qquad \textbf{(D)}\ 14\qquad \textbf{(E)}\ 15$


	\item (\textit{AHSME 1988}) In an experiment, a scientific constant $C$ is determined to be $2.43865$ with an error of at most $\pm 0.00312$. 
The experimenter wishes to announce a value for $C$ in which every digit is significant. 
That is, whatever $C$ is, the announced value must be the correct result when $C$ is rounded to that number of digits. 
The most accurate value the experimenter can announce for $C$ is


 $\textbf{(A)}\ 2\qquad \textbf{(B)}\ 2.4\qquad \textbf{(C)}\ 2.43\qquad \textbf{(D)}\ 2.44\qquad \textbf{(E)}\ 2.439$


	\item (\textit{AHSME 1988}) 
                    \begin{figure}[H]
                        \centering
                        \includegraphics[width=0.5\textwidth]{Problem/AHSME_I_1988_11_Problem_Figure_1.png}
                    \end{figure}
                    


 On each horizontal line in the figure below, the five large dots indicate the populations of cities $A, B, C, D$ and $E$ in the year indicated. 
Which city had the greatest percentage increase in population from $1970$ to $1980$?


 $\textbf{(A)}\ A\qquad \textbf{(B)}\ B\qquad \textbf{(C)}\ C\qquad \textbf{(D)}\ D\qquad \textbf{(E)}\ E$


	\item (\textit{AHSME 1988}) Each integer $1$ through $9$ is written on a separate slip of paper and all nine slips are put into a hat. 
Jack picks one of these slips at random and puts it back. Then Jill picks a slip at random. 
Which digit is most likely to be the units digit of the sum of Jack's integer and Jill's integer?


 $\textbf{(A)}\ 0\qquad \textbf{(B)}\ 1\qquad \textbf{(C)}\ 8\qquad \textbf{(D)}\ 9\qquad \textbf{(E)}\ \text{each digit is equally likely}$


	\item (\textit{AHSME 1988}) If $\sin(x)= 3\cos(x)$ then what is $\sin(x) \cdot \cos(x)$?


 $\textbf{(A)}\ \frac{1}{6}\qquad \textbf{(B)}\ \frac{1}{5}\qquad \textbf{(C)}\ \frac{2}{9}\qquad \textbf{(D)}\ \frac{1}{4}\qquad \textbf{(E)}\ \frac{3}{10}$


	\item (\textit{AHSME 1988}) For any real number a and positive integer k, define


 ${a \choose k} = \frac{a(a-1)(a-2)\cdots(a-(k-1))}{k(k-1)(k-2)\cdots(2)(1)}$


 What is


 ${-\frac{1}{2} \choose 100} \div {\frac{1}{2} \choose 100}$?


 $\textbf{(A)}\ -199\qquad \textbf{(B)}\ -197\qquad \textbf{(C)}\ -1\qquad \textbf{(D)}\ 197\qquad \textbf{(E)}\ 199$


	\item (\textit{AHSME 1988}) If $a$ and $b$ are integers such that $x^2 - x - 1$ is a factor of $ax^3 + bx^2 + 1$, then $b$ is


 $\textbf{(A)}\ -2\qquad \textbf{(B)}\ -1\qquad \textbf{(C)}\ 0\qquad \textbf{(D)}\ 1\qquad \textbf{(E)}\ 2$


	\item (\textit{AHSME 1988}) 
                    \begin{figure}[H]
                        \centering
                        \includegraphics[width=0.5\textwidth]{Problem/AHSME_I_1988_16_Problem_Figure_1.png}
                    \end{figure}
                    
$ABC$ and $A'B'C'$ are equilateral triangles with parallel sides and the same center, 
as in the figure. The distance between side $BC$ and side $B'C'$ is $\frac{1}{6}$ the altitude of $\triangle ABC$. 
The ratio of the area of $\triangle A'B'C'$ to the area of $\triangle ABC$ is


 $\textbf{(A)}\ \frac{1}{36}\qquad \textbf{(B)}\ \frac{1}{6}\qquad \textbf{(C)}\ \frac{1}{4}\qquad \textbf{(D)}\ \frac{\sqrt{3}}{4}\qquad \textbf{(E)}\ \frac{9+8\sqrt{3}}{36}$


	\item (\textit{AHSME 1988}) If $|x| + x + y = 10$ and $x + |y| - y = 12$, find $x + y$
$\textbf{(A)}\ -2\qquad \textbf{(B)}\ 2\qquad \textbf{(C)}\ \frac{18}{5}\qquad \textbf{(D)}\ \frac{22}{3}\qquad \textbf{(E)}\ 22$


	\item (\textit{AHSME 1988}) At the end of a professional bowling tournament, the top 5 bowlers have a playoff. 
First \#5 bowls \#4. The loser receives $5$th prize and the winner bowls \#3 in another game. 
The loser of this game receives $4$th prize and the winner bowls \#2. 
The loser of this game receives $3$rd prize and the winner bowls \#1. 
The winner of this game gets 1st prize and the loser gets 2nd prize. 
In how many orders can bowlers \#1 through \#5 receive the prizes?


 $\textbf{(A)}\ 10\qquad \textbf{(B)}\ 16\qquad \textbf{(C)}\ 24\qquad \textbf{(D)}\ 120\qquad \textbf{(E)}\ \text{none of these}$


	\item (\textit{AHSME 1988}) Simplify


 $\frac{bx(a^2x^2 + 2a^2y^2 + b^2y^2) + ay(a^2x^2 + 2b^2x^2 + b^2y^2)}{bx + ay}$
$\textbf{(A)}\ a^2x^2 + b^2y^2\qquad \textbf{(B)}\ (ax + by)^2\qquad \textbf{(C)}\ (ax + by)(bx + ay)\qquad\\ \textbf{(D)}\ 2(a^2x^2+b^2y^2)\qquad \textbf{(E)}\ (bx+ay)^2$


	\item (\textit{AHSME 1988}) In one of the adjoining figures a square of side $2$ is dissected into four pieces so that $E$ and $F$ are the midpoints
of opposite sides and $AG$ is perpendicular to $BF$. These four pieces can then be reassembled into a rectangle as shown 
in the second figure. The ratio of height to base, $XY / YZ$, in this rectangle is  


 
                    \begin{figure}[H]
                        \centering
                        \includegraphics[width=0.5\textwidth]{Problem/AHSME_I_1988_20_Problem_Figure_1.png}
                    \end{figure}
                    
$\textbf{(A)}\ 4\qquad \textbf{(B)}\ 1+2\sqrt{3}\qquad \textbf{(C)}\ 2\sqrt{5}\qquad \textbf{(D)}\ \frac{8+4\sqrt{3}}{3}\qquad \textbf{(E)}\ 5$


	\item (\textit{AHSME 1988}) The complex number $z$ satisfies $z + |z| = 2 + 8i$. What is $|z|^{2}$? Note: if $z = a + bi$, then $|z| = \sqrt{a^{2} + b^{2}}$.


 $\textbf{(A)}\ 68\qquad \textbf{(B)}\ 100\qquad \textbf{(C)}\ 169\qquad \textbf{(D)}\ 208\qquad \textbf{(E)}\ 289$


	\item (\textit{AHSME 1988}) For how many integers $x$ does a triangle with side lengths $10, 24$ and $x$ have all its angles acute?


 $\textbf{(A)}\ 4\qquad \textbf{(B)}\ 5\qquad \textbf{(C)}\ 6\qquad \textbf{(D)}\ 7\qquad \textbf{(E)}\ \text{more than } 7$


	\item (\textit{AHSME 1988}) The six edges of a tetrahedron $ABCD$ measure $7, 13, 18, 27, 36$ and $41$ units. If the length of edge $AB$ is $41$, then the length of edge $CD$ is


 $\textbf{(A)}\ 7\qquad \textbf{(B)}\ 13\qquad \textbf{(C)}\ 18\qquad \textbf{(D)}\ 27\qquad \textbf{(E)}\ 36$


	\item (\textit{AHSME 1988}) An isosceles trapezoid is circumscribed around a circle. The longer base of the trapezoid is $16$, 
and one of the base angles is $\arcsin(.8)$. Find the area of the trapezoid.


 $\textbf{(A)}\ 72\qquad \textbf{(B)}\ 75\qquad \textbf{(C)}\ 80\qquad \textbf{(D)}\ 90\qquad \textbf{(E)}\ \text{not uniquely determined}$


	\item (\textit{AHSME 1988}) $X, Y$ and $Z$ are pairwise disjoint sets of people. The average ages of people in the sets 
$X, Y, Z, X \cup Y, X \cup Z$ and $Y \cup Z$ are $37, 23, 41, 29, 39.5$ and $33$ respectively. 
Find the average age of the people in set $X \cup Y \cup Z$.


 $\textbf{(A)}\ 33\qquad \textbf{(B)}\ 33.5\qquad \textbf{(C)}\ 33.6\overline{6}\qquad \textbf{(D)}\ 33.83\overline{3}\qquad \textbf{(E)}\ 34$


	\item (\textit{AHSME 1988}) Suppose that $p$ and $q$ are positive numbers for which


 $\log_{9}(p) = \log_{12}(q) = \log_{16}(p+q)$


 What is the value of $\frac{q}{p}$?


 $\textbf{(A)}\ \frac{4}{3}\qquad \textbf{(B)}\ \frac{1+\sqrt{3}}{2}\qquad \textbf{(C)}\ \frac{8}{5}\qquad \textbf{(D)}\ \frac{1+\sqrt{5}}{2}\qquad \textbf{(E)}\ \frac{16}{9}$


	\item (\textit{AHSME 1988}) In the figure, $AB \perp BC, BC \perp CD$, and $BC$ is tangent to the circle with center $O$ and diameter $AD$. 
In which one of the following cases is the area of $ABCD$ an integer?


 
                    \begin{figure}[H]
                        \centering
                        \includegraphics[width=0.5\textwidth]{Problem/AHSME_I_1988_27_Problem_Figure_1.png}
                    \end{figure}
                    
$\textbf{(A)}\ AB=3, CD=1\qquad \textbf{(B)}\ AB=5, CD=2\qquad \textbf{(C)}\ AB=7, CD=3\qquad\\ \textbf{(D)}\ AB=9, CD=4\qquad \textbf{(E)}\ AB=11, CD=5$


	\item (\textit{AHSME 1988}) An unfair coin has probability $p$ of coming up heads on a single toss. 
Let $w$ be the probability that, in $5$ independent toss of this coin, 
heads come up exactly $3$ times. If $w = 144 / 625$, then 


 $\textbf{(A)}\ p\text{ must be }\frac{2}{5}\qquad \textbf{(B)}\ p\text{ must be }\frac{3}{5}\qquad\\ \textbf{(C)}\ p\text{ must be greater than }\frac{3}{5}\qquad \textbf{(D)}\ p\text{ is not uniquely determined}\qquad\\ \textbf{(E)}\ \text{there is no value of } p \text{ for which }w =\frac{144}{625}$


	\item (\textit{AHSME 1988}) You plot weight $(y)$ against height $(x)$ for three of your friends and obtain the points 
$(x_{1},y_{1}), (x_{2},y_{2}), (x_{3},y_{3})$. If $x_{1} < x_{2} < x_{3}$ and $x_{3} - x_{2} = x_{2} - x_{1}$, 
which of the following is necessarily the slope of the line which best fits the data? 
"Best fits" means that the sum of the squares of the vertical distances from the data points to the line is smaller than for any other line.


 $\textbf{(A)}\ \frac{y_{3}-y_{1}}{x_{3}-x_{1}}\qquad \textbf{(B)}\ \frac{(y_{2}-y_{1})-(y_{3}-y_{2})}{x_{3}-x_{1}}\qquad\\ \textbf{(C)}\ \frac{2y_{3}-y_{1}-y_{2}}{2x_{3}-x_{1}-x_{2}}\qquad \textbf{(D)}\ \frac{y_{2}-y_{1}}{x_{2}-x_{1}}+\frac{y_{3}-y_{2}}{x_{3}-x_{2}}\qquad\\ \textbf{(E)}\ \text{none of these}$


	\item (\textit{AHSME 1988}) Let $f(x) = 4x - x^{2}$. Give $x_{0}$, consider the sequence defined by $x_{n} = f(x_{n-1})$ for all $n \ge 1$. 
For how many real numbers $x_{0}$ will the sequence $x_{0}, x_{1}, x_{2}, \ldots$ take on only a finite number of different values?


 $\textbf{(A)}\ \text{0}\qquad \textbf{(B)}\ \text{1 or 2}\qquad \textbf{(C)}\ \text{3, 4, 5 or 6}\qquad\\ \textbf{(D)}\ \text{more than 6 but finitely many}\qquad\\ \textbf{(E) }\infty$


\end{enumerate}